% ID: SEL-D0038-OB-E01
% ESTADO: validada. S039, cotejada con examen y criterios.
\textbf{a)} Sean $x,y$ los números de trajes y abrigos.
Las restricciones son $x,y\ge0$, $x+2y\le160$ y $x+y\le120$.
Los precios permiten maximizar el ingreso $I=250x+350y$,
al que el enunciado llama beneficio.

En los ejes se obtienen $(120,0)$ y $(0,80)$.
Las fronteras de los tejidos se cruzan en $(80,40)$:
al restar sus ecuaciones queda $y=40$.
\begin{center}
\setlength{\unitlength}{7mm}
\begin{picture}(11,7)(-2,-1)
\qbezier(0,0)(4.25,0)(8.5,0)\put(8.7,0){$x$}
\qbezier(0,0)(0,2.85)(0,5.7)\put(0,5.8){$y$}
\qbezier(0,5.33333)(2.66667,4)(5.33333,2.66667)
\qbezier(5.33333,2.66667)(6.66667,1.33333)(8,0)
\put(-0.1,5.33){\makebox(0,0)[r]{$(0,80)$}}
\put(5.5,2.9){$(80,40)$}
\put(8,-0.3){\makebox(0,0)[t]{$(120,0)$}}
\end{picture}
\end{center}
En $(0,0)$, $(120,0)$, $(80,40)$ y $(0,80)$ los ingresos son
0, 30000, 34000 y 28000 euros. Por linealidad, el óptimo se
alcanza en el mejor vértice, que tiene coordenadas enteras.
Son \textbf{80 trajes y 40 abrigos}, con $\boxed{34000}$ euros.

\textbf{b)} Con $(60,50)$ se usan $60+2(50)=160$ m de pana
y $2(60+50)=220$ m de lana. Es factible, pero los ingresos
son $250(60)+350(50)=32500$ euros, inferiores al máximo.
